% Appendix: construction of the outer scalar for p different from 1.
\providecommand{\eps}{\varepsilon}
\providecommand{\Log}{\operatorname{Log}}
\providecommand{\re}{\operatorname{Re}}
\providecommand{\im}{\operatorname{Im}}
\providecommand{\thlip}{\vartheta_{\mathrm{lip}}}
\providecommand{\thseam}{\vartheta_{\mathrm{seam}}}
\providecommand{\thw}{\vartheta_{\mathrm{w}}}
\providecommand{\Ob}{\Omega_{\mathrm{b}}}
\providecommand{\Sin}{S_{\mathrm{in}}}
\providecommand{\Sout}{S_{\mathrm{out}}}
\providecommand{\Zn}{\mathcal{Z}_n}
\providecommand{\Sseam}{\Sigma^{\mathrm{s}}}

\section{\texorpdfstring{Construction of the outer scalar for $p\ne1$}{Construction of the outer scalar for p different from 1}}\label{app:outer}

This appendix proves Lemma~\ref{lem:outer} for $p>1$ (\S\ref{sec:outer-pgt1}) and for $0<p<1$
(\S\ref{sec:outer-plt1}). We use the notation of Section~\ref{sec:outer}: Binet's function $j$, $\tilde D=e^{-j(t)}$,
$D=\kappa_\pm\tilde D$, $K=(4c-z)^{-1/2}$, and the bounds~\eqref{eq:binet-bound}, \eqref{eq:binet-small}. In both cases
the difficulty is the cut of $t(z)$ on $(-\infty,0)$, where $D$ has a jump that is not of the form required by (O1).
For $p>1$ the jump is removed exactly by a Cauchy integral; for $p<1$, $D$ is used only in an inner sector and glued
to an elementary function outside.

\subsection{\texorpdfstring{The case $p>1$: the Binet-cut scalar}{The case p>1: the Binet-cut scalar}}\label{sec:outer-pgt1}
For $p>1$, $|\arg t|<\pi/p<\pi$ on the slit plane, so $\Gamma(t)$ has no poles there, and the jump of $D$ on
$(-\infty,0)$ is a bounded phase $e^{i\mathfrak b}$.

\begin{lemma}[Binet-cut scalar]\label{lem:binetcut}
Let $p>1$.
\begin{enumerate}
\item[(i)] $\tilde D$ is analytic and zero-free on $\C\setminus(-\infty,0]$, and $D$ is analytic and zero-free in
  $\C_\pm$, with continuous boundary values on $\R\setminus(\{X_k\}\cup\{0\})$. On $(-\infty,0)$, $D_-=\overline{D_+}$.
\item[(ii)] Let $\mathfrak b(y):=2\im\bigl(-j(t_+(y))-\Log(1-e^{-2i\theta_+(y)})\bigr)$ for $y<0$, where $t_+,\theta_+$
  are boundary values from $\C_+$. Then $D_+/D_-=e^{i\mathfrak b}$ on $(-\infty,0)$; $\mathfrak b$ is real-analytic,
  $|\mathfrak b(y)|\le C\min(1,|y|^{1/p})$ with $C$ independent of $n$, $|\mathfrak b'(y)|\le C_n|y|^{1/p-1}$ for
  $-1\le y<0$, $\mathfrak b'$ is bounded on $(-\infty,-\frac12]$, and $\mathfrak b(y)\to\pi/p-\pi$ as $y\to-\infty$, with
  $\int_{-\infty}^{-1}|\mathfrak b(y)-\pi/p+\pi|\,|y|^{-1}\diff y<\infty$.
\item[(iii)] The function
  \begin{equation}\label{eq:B}
    B(z):=\exp\Bigl(-\frac1{2\pi}\int_{-\infty}^0\mathfrak b(y)\Bigl[\frac1{y-z}-\frac1{y-1}\Bigr]\diff y\Bigr)
  \end{equation}
  is analytic and zero-free on $\C\setminus(-\infty,0]$, positive on $(0,\infty)$, has continuous boundary values on
  $(-\infty,0)$ with $B_+=B_-e^{-i\mathfrak b}$, and $DB$ extends analytically across $(-\infty,0)$. For $z\in\Zn$, with
  $C$ independent of $n$,
  \begin{equation}\label{eq:B-bounds}
  \begin{gathered}
    |\log B(z)-\log B(0)|\le C|z|^{1/p},\quad |(\log B)'(z)|\le C|z|^{1/p-1}\quad(|z|\le1),\\
    |\log B(z)|\le C(1+\log(1+|z|)).
  \end{gathered}
  \end{equation}
\item[(iv)] $F:=DBK$ satisfies (O1)--(O5), with $\tilde F=\tilde DBK$ and $H_n=(\tilde DB)^{-2}$, $H_n(0^+)=B(0)^{-2}$.
\end{enumerate}
\end{lemma}

\begin{proof}
Part~(i) is the absence of poles of $\Gamma$ in the sector; (ii) estimates the jump $\mathfrak b$ by Stirling's formula
for large $|t|$ and by~\eqref{eq:binet-small} for small $|t|$; (iii) is the Sokhotski--Plemelj formula together with a
splitting of the Cauchy integral at $|y|=|z|$; (iv) checks (O1)--(O5) one by one.

(i) On the slit plane $|\arg t|<\pi/p<\pi$, so $\Gamma(t)$ has neither zeros nor poles and $\tilde D$ is analytic and
zero-free. By~\eqref{eq:imtheta}, $|e^{\mp2i\theta}|<1$ in $\C_\pm$ (here $\pi p>\pi$), so $D$ is analytic and
zero-free in $\C_\pm$. On $(-\infty,0)$, $t_-=\overline{t_+}$ and $\Gamma$, $j$ are real on the real axis, so
$D_-=\overline{D_+}$.

(ii) Since $|e^{-2i\theta_+}|<1$ on $(-\infty,0)$, $\log D_+:=-j(t_+)-\Log(1-e^{-2i\theta_+})$ is a continuous
logarithm, real-analytic in $y$, and $D_+/D_-=e^{2i\im\log D_+}=e^{i\mathfrak b}$. With $|t_+|=n|y|^{-1/p}$ and
$\arg t_+=-\pi/p$: for $|t_+|\ge1$, \eqref{eq:binet-bound} and
$|\Log(1-e^{-2i\theta_+})|\le2\exp(-2\pi\sin(\pi/p)|t_+|)$ give $|\mathfrak b|\le C/|t_+|=C|y|^{1/p}/n$. For $|t_+|\le1$
the expansion~\eqref{eq:binet-small} and $\Log(1-e^{-2\pi it})=\log(2\pi|t|)+i(\frac\pi2+\arg t)+O(t)$ (note
$\frac\pi2+\arg t_+\in(-\frac\pi2,\frac\pi2)$) show that $\mathfrak b$ is bounded and
\[
  \mathfrak b(y)=2\bigl(\tfrac12\arg t_+-\tfrac\pi2-\arg t_+\bigr)+O(|t_+|\log(2/|t_+|))=\frac\pi p-\pi+O\bigl(n|y|^{-1/p}\log|y|\bigr).
\]
This gives the bounds on $\mathfrak b$ and the integrability. For the derivative, $\diff t_+/\diff y=-t_+/(py)$, and
$\frac{\diff}{\diff t}\bigl(j(t)+\Log(1-e^{-2\pi it})\bigr)$ is $O(|t|^{-2})$ for $|t|\ge1$ and $O(|t|^{-1})$ for
$|t|\le1$; multiplied by $|t|/|y|$ this gives $|\mathfrak b'(y)|\le C|y|^{1/p-1}/n$ for $|y|\le n^p$ and
$|\mathfrak b'(y)|\le C/|y|$ otherwise.

(iii) The integral converges absolutely, since $\mathfrak b$ is bounded and the bracket is $(z-1)/((y-z)(y-1))$. By the
Sokhotski--Plemelj formula (the density is real-analytic on $(-\infty,0)$), $\log B$ has continuous boundary values
and $(\log B)_+-(\log B)_-=-i\mathfrak b$. So $(DB)_+=(DB)_-$ on $(-\infty,0)$, and $DB$, being continuous across, is
analytic there by Morera's theorem. $B$ is positive on $(0,\infty)$ because the integrand is real there.

For~\eqref{eq:B-bounds}, let $z\in\Zn$. Then $|y-z|^2=y^2+|z|^2+2|y||z|\cos(\arg z)\ge\frac12(|y|+|z|)^2$ for $y<0$,
since $|\arg z|<\thw<\pi/2$. We have $\log B(z)-\log B(0)=-\frac z{2\pi}\int_{-\infty}^0\frac{\mathfrak b(y)\diff y}{y(y-z)}$, so
by~(ii), for $|z|\le1$,
\[
  |\log B(z)-\log B(0)|\le C|z|\int_0^\infty\frac{\min(1,u^{1/p})}{u(u+|z|)}\diff u\le
  C|z|\Bigl(\int_0^1\frac{u^{1/p-1}\diff u}{u+|z|}+1\Bigr)\le C|z|^{1/p},
\]
because $1/p-1\in(-1,0)$. The same splitting applied to $(\log B)'(z)=-\frac1{2\pi}\int\mathfrak b(y)(y-z)^{-2}\diff y$
gives $C|z|^{1/p-1}$, and applied to the original integral gives $C(1+\log(1+|z|))$.

(iv) \emph{(O1).} $F$ is analytic and zero-free off $[0,\infty)$ by~(i), (iii) and the choice of $K$. On
$(4c,\infty)$ off the nodes, $F_+/F_-=(D_+/D_-)(K_+/K_-)=(-e^{2i\theta})(-1)$, since $B$ is analytic there.

\emph{(O2).} $\tilde F=F/\kappa_\pm=\tilde DBK$ is analytic and zero-free on $\C\setminus\bigl((-\infty,0]\cup[4c,\infty)\bigr)$,
and $\tilde F_+=-\tilde F_-$ on $(4c,\infty)$ because of $K$. Thus $H_n=1/((4c-z)\tilde F^2)=(\tilde DB)^{-2}$, which is
positive on $(0,\infty)$, and $H_n(0^+)=B(0)^{-2}$ because $\tilde D(z)\to1$ as $z\to0$ in $\Zn$ by~\eqref{eq:binet-bound}.

\emph{(O3), (O4).} $\log H_n-\log H_n(0^+)=2j(t)-2(\log B-\log B(0))$. For $z\in\Zn$, $|z|\le1$, we have $|t|\ge n$ and
$|\arg t|\le\thw/p$, so $|j(t)|\le C|z|^{1/p}/n$ and $|\frac{\diff}{\diff z}j(t)|=|j'(t)||t|/(p|z|)\le C|z|^{1/p-1}/n$.
Together with~\eqref{eq:B-bounds} and $\nu_p\le1/p$, this is (O3). For $|z|\ge1$ in $\Zn$ we have $|t|\ge e^{-c_0}$, so
$j(t)$ is bounded (by~\eqref{eq:binet-bound} for $|t|\ge1$ and by continuity on the compact set
$\{e^{-c_0}\le|t|\le1,\ |\arg t|\le\thw/p\}$), and~\eqref{eq:B-bounds} gives (O4).

\emph{(O5).} Fix $n$. Near $0$ and outside the lens, $|\arg z|\ge\frac34\thlip$, so $\kappa_\pm\to1$ by~\eqref{eq:imtheta}
and $\tilde D\to1$; $K$ is bounded; and $\log B$ is bounded on $U(0,\delta)\setminus(-\delta,0]$ uniformly up to the cut.
For the last point split $\int\mathfrak b(y)/(y-z)\diff y$ into the window $|y-\re z|\le|\re z|/2$ and the rest. On the
window write $\mathfrak b(y)=\mathfrak b(\re z)+O(|y-\re z|\,|\re z|^{1/p-1})$ by~(ii); the constant part integrates to
$\mathfrak b(\re z)\Log(\cdot)$ with a bounded logarithm (the end points of the window are at comparable distances from
$z$), and the rest is bounded; outside the window $|y-z|\ge(|y|+|z|)/C$ and $|\mathfrak b(y)|\le C|y|^{1/p}$ give a
bounded contribution. In the lens, $\tilde F=\tilde DBK$ is bounded above and below by the same facts.

Near $\infty$, the same window argument on $(-\infty,-\frac12]$, where $\mathfrak b'$ is bounded, gives
$|\log B(z)|\le C_n(1+\log|z|)$ for $|z|\ge2$ off the cut. For $|z|\ge n^p+\delta'$ we have $|t|<1$, and $t$ stays away from
$1$, so $D$ is bounded above and below by a power of $|z|$; hence $|F|^{\pm1}\le C_n|z|^{N_n}$. Finally, along $z=iy$,
$y\to\pm\infty$: by~\eqref{eq:binet-small}, $|D(iy)|=|t|^{-1/2}(2\pi)^{-1/2}(1+o(1))\asymp|y|^{1/(2p)}$; $|K(iy)|\asymp|y|^{-1/2}$;
and splitting $\mathfrak b=(\pi/p-\pi)\mathbf 1_{(-\infty,-1)}+\mathfrak e$, the first part contributes
$-\frac{\pi/p-\pi}{2\pi}\Log\frac{z+1}{2}$ to $\log B$ (an explicit integral), while the second contributes a quantity with
a finite limit, by dominated convergence and the integrability in~(ii). So $|B(iy)|\asymp|y|^{\frac12-\frac1{2p}}$, and
$|F(iy)|\asymp1$.
\end{proof}

\subsection{\texorpdfstring{The case $0<p<1$: seams}{The case 0<p<1: seams}}\label{sec:outer-plt1}
For $p<1$ the poles $t=-1,-2,\dots$ of $\Gamma$ lie on the rays $\arg z=\pm\pi p$, inside the slit plane, and $D$ has
no useful continuation to the negative axis. We use $D$ only in the inner sector, where $|\arg t|<\pi/2$, and an
elementary function outside; a Cauchy integral $\chi$ over the two separating rays glues them together exactly. Put
\begin{gather*}
  \thseam:=\tfrac{\pi p}2,\qquad \Sin:=\{z\ne0:\ |\arg z|<\thseam\},\qquad \Sout:=\C\setminus(\overline{\Sin}\cup\{0\}),\\
  \Sseam_\pm:=\{re^{\pm i\thseam}:\ r>0\},
\end{gather*}
both seams oriented away from $0$, and
\[
  \varpi:=\tfrac12-\tfrac1{2p}<0,\qquad B_{\mathrm{in}}(z):=(1+z)^{\varpi},\qquad B_{\mathrm{out}}(z):=\Bigl(1-\frac z{2\pi n}\Bigr)^{1/2},
\]
with principal powers; the cut $(-\infty,-1]$ of $B_{\mathrm{in}}$ lies in $\Sout$, and the cut $[2\pi n,\infty)$ of
$B_{\mathrm{out}}$ lies in $\Sin$. The lens, the saturated regions and the top segment lie in $\Zn\subset\Sin$, since
$\thw<\thseam$ by~\eqref{eq:angles}.

\begin{lemma}[Seam scalar]\label{lem:seam}
Let $0<p<1$.
\begin{enumerate}
\item[(i)] On $\Sin$, $|\arg t|<\pi/2$; $\tilde D=e^{-j(t)}$ is analytic and zero-free there, and $D=\kappa_\pm\tilde D$ is
  analytic and zero-free in $\Sin\cap\C_\pm$.
\item[(ii)] On $\Sseam_\pm$ put, with principal logarithms,
  \begin{equation}\label{eq:Lpm}
    \chi_\pm(\zeta):=-j(t(\zeta))-\Log\bigl(1-e^{\mp2\pi it(\zeta)}\bigr)+\varpi\Log(1+\zeta)-\tfrac12\Log\Bigl(1-\frac\zeta{2\pi n}\Bigr).
  \end{equation}
  Then $\chi_\pm$ is real-analytic on the open ray, $e^{\chi_\pm}=DB_{\mathrm{in}}/B_{\mathrm{out}}$ there,
  $\chi_\pm(\zeta)\to0$ as $\zeta\to0$, and $\chi_\pm(\zeta)=O\bigl(n|\zeta|^{-1/p}\log|\zeta|+n|\zeta|^{-1}\bigr)$ as
  $\zeta\to\infty$. In particular $\chi_\pm$ has no winding: it is the continuous logarithm of
  $DB_{\mathrm{in}}/B_{\mathrm{out}}$ that vanishes at both ends. Moreover $\chi_-(\bar\zeta)=\overline{\chi_+(\zeta)}$.
\item[(iii)] The \emph{seam scalar}
  \begin{equation}\label{eq:chi}
    \chi(z):=\frac1{2\pi i}\Bigl[\int_{\Sseam_+}\frac{\chi_+(\zeta)\diff\zeta}{\zeta-z}-\int_{\Sseam_-}\frac{\chi_-(\zeta)\diff\zeta}{\zeta-z}\Bigr]
  \end{equation}
  is analytic off $\Sseam_+\cup\Sseam_-$, real on $\R\setminus\{0\}$, tends to $0$ at $\infty$, has continuous boundary
  values on the open seams, and $\chi_{\mathrm{out}}-\chi_{\mathrm{in}}=\chi_\pm$ on $\Sseam_\pm$, where ``out'' and
  ``in'' refer to the boundary values from $\Sout$ and $\Sin$.
\item[(iv)] The function
  \begin{equation}\label{eq:F-seam}
    F:=DB_{\mathrm{in}}Ke^{\chi}\ \text{in }\Sin,\qquad F:=B_{\mathrm{out}}Ke^{\chi}\ \text{in }\Sout,
  \end{equation}
  extends analytically across $\Sseam_\pm$ and satisfies (O1)--(O5), with $\tilde F=\tilde DB_{\mathrm{in}}Ke^\chi$ and
  $H_n=(\tilde DB_{\mathrm{in}}e^\chi)^{-2}$, $H_n(0^+)=e^{-2\chi(0)}$.
\end{enumerate}
\end{lemma}

\begin{proof}
(i) $|\arg t|=|\arg z|/p<\pi/2$ on $\Sin$, so $\Gamma(t)$ has no poles or zeros there; and $\pm\im\theta<0$ in
$\Sin\cap\C_\pm$ by~\eqref{eq:imtheta}.

(ii) Each term of~\eqref{eq:Lpm} is continuous on the open ray $\Sseam_+$: $j$ is analytic off $(-\infty,0]$;
$\arg t=-\pi/2$, so $e^{-2\pi it}=e^{-2\pi|t|}\in(0,1)$; $1+\zeta\notin(-\infty,0]$ since $\arg\zeta\in(0,\pi)$; and
$1-\zeta/(2\pi n)$ has non-zero imaginary part. The exponential of the sum is $\tilde D\kappa_+B_{\mathrm{in}}/B_{\mathrm{out}}$.
As $\zeta\to0$, $|t|\to\infty$ and each term tends to $0$ (by~\eqref{eq:binet-bound} for $j$). As $\zeta\to\infty$,
$t\to0$ and, by~\eqref{eq:binet-small}, using $\arg t=-\pi/2$, $\arg\zeta=\thseam$ and $\Log t=\log n-\frac1p\Log\zeta$:
\begin{align*}
  -j(t)&=\tfrac12\Log t+\tfrac12\log(2\pi)+O(|t|\log(2/|t|)),\\
  -\Log(1-e^{-2\pi it})&=-\log(2\pi)-\Log t-\tfrac{i\pi}2+O(|t|),\\
  \varpi\Log(1+\zeta)&=\varpi\Log\zeta+O(|\zeta|^{-1}),\\
  -\tfrac12\Log\bigl(1-\tfrac\zeta{2\pi n}\bigr)&=-\tfrac12\Log\zeta+\tfrac12\log(2\pi n)+\tfrac{i\pi}2+O(n|\zeta|^{-1});
\end{align*}
the second line uses $1-e^{-2\pi it}=2\pi it(1+O(t))$ and $\arg(it)=0$, the last uses
$\arg(-\zeta/(2\pi n))=\thseam-\pi\in(-\pi,0)$. The sum of the main terms is
$-\frac12\Log t+\frac12\log n+(\varpi-\frac12)\Log\zeta=\bigl(\frac1{2p}+\varpi-\frac12\bigr)\Log\zeta=0$. The symmetry
follows from $\overline{t(\bar\zeta)}=t(\zeta)$ and the reality of $j$.

(iii) By~(ii), $\chi_\pm(\zeta)/\zeta$ is integrable on $\Sseam_\pm$ near $0$ and at $\infty$ (since $1/p>1$), so the
integrals converge and $\chi\to0$ at $\infty$. The Sokhotski--Plemelj formula gives the jumps: $\Sseam_+$ is oriented
outwards with $\Sout$ on its left, so $\chi_{\mathrm{out}}-\chi_{\mathrm{in}}=\chi_+$; on $\Sseam_-$ the left side is
$\Sin$ and the density carries a minus sign, so again $\chi_{\mathrm{out}}-\chi_{\mathrm{in}}=\chi_-$. Reality on $\R$
follows from the symmetry in~(ii).

(iv) On $\Sseam_\pm$, $F_{\mathrm{in}}/F_{\mathrm{out}}=DB_{\mathrm{in}}e^{\chi_{\mathrm{in}}}/(B_{\mathrm{out}}e^{\chi_{\mathrm{out}}})
=e^{\chi_\pm}e^{-\chi_\pm}=1$, so $F$ is continuous, hence analytic, across the seams.

\emph{(O1).} In $\Sout$ every factor is analytic and zero-free, and $(-\infty,0)\subset\Sout$. In $\Sin\setminus[0,\infty)$
the same holds by~(i). On $(4c,\infty)$, $B_{\mathrm{in}}$ and $\chi$ are analytic, so $F_+/F_-=(-e^{2i\theta})(-1)$.

\emph{(O2).} $\tilde F=\tilde DB_{\mathrm{in}}Ke^\chi$ on $\Sin\supset\Zn$, analytic and zero-free off $[4c,\infty)$, with
$\tilde F_+=-\tilde F_-$ there. $H_n=(\tilde DB_{\mathrm{in}}e^{\chi})^{-2}>0$ on $(0,\infty)$ by~(iii), and
$H_n(0^+)=e^{-2\chi(0)}$, where $\chi(0):=\lim_{z\to0}\chi(z)$ exists by Lemma~\ref{lem:corner}.

\emph{(O3), (O4).} Points of $\Zn$ satisfy $|\arg z|<\thw$, hence $\mathrm{dist}(z,\Sseam_+\cup\Sseam_-)\ge\sin(\thseam-\thw)|z|$
(note $0<\thseam-\thw<\pi/2$). So Lemma~\ref{lem:corner}(ii) applies with $\tau_0=\sin(\thseam-\thw)$. For $|z|\le1$,
\begin{align*}
  |\log H_n-\log H_n(0^+)|&\le2\bigl(|j(t)|+|\varpi||\Log(1+z)|+|\chi-\chi(0)|\bigr)\\
  &\le C\bigl(|z|^{1/p}/n+|z|+|z|(1+\log\tfrac1{|z|})\bigr)\le C|z|^{1/2},
\end{align*}
and similarly $|(\log H_n)'|\le C(|z|^{1/p-1}/n+1+\log\frac1{|z|})\le C|z|^{-1/2}$. Since $\nu_p=\frac12$, this is (O3). For
$|z|\ge1$ in $\Zn$, $j(t)$ is bounded, $|\Log(1+z)|\le C(1+\log|z|)$, and $|\chi-\chi(0)|\le C(1+\log|z|)^2$ by
Lemma~\ref{lem:corner}(ii); this is (O4).

\emph{(O5).} Fix $n$. Near $0$: in $\Sin$ outside the lens, $\kappa_\pm\to1$ (as $\frac34\thlip/p\le|\arg t|<\frac\pi2$ there, so that
$\sin|\arg t|\ge\sin(\frac34\thlip/p)>0$ in~\eqref{eq:imtheta}) and $\tilde D,B_{\mathrm{in}}\to1$; in $\Sout$, $B_{\mathrm{out}}\to1$; $K$ is bounded; and $\chi$ is
bounded near $0$ uniformly up to the seams by Lemma~\ref{lem:corner}(iii). In the lens, $\tilde F$ is bounded above and
below by the same facts. Near $\infty$: in $\Sout$ (which contains the imaginary axis, as $\thseam<\pi/2$),
$F=B_{\mathrm{out}}Ke^\chi\to(2\pi n)^{-1/2}$, because $B_{\mathrm{out}}(z)K(z)=\bigl(\frac{-z}{2\pi n}\bigr)^{1/2}(-z)^{-1/2}(1+o(1))$
with consistent principal branches in $\Sout$. In $\Sin\cap\{|z|\ge n^p+\delta'\}$, $|t|<1$ and $t$ stays away from $1$, so
$D$, $B_{\mathrm{in}}$, $K$ and $e^\chi$ are bounded above and below by powers of $|z|$.
\end{proof}

The seams carry no weight: $\chi_\pm$ and $\chi$ are built from $j$, $e^{\mp2\pi it}$, $B_{\mathrm{in}}$ and
$B_{\mathrm{out}}$ only, and $F$ has no jump across them. Neither $V$ nor $Q$ has a jump on the seams, and the
comparison parametrix of Section~\ref{sec:hardedge} has none either, since it is built from $D_W$. The only trace of the
seams in the hard-edge analysis is the factor $e^{-2\chi}$ in $H_n$, which is controlled by the following lemma. The
constant $\chi(0)$ depends on $n$, but no bound on it is needed: it enters only through $H_n(0^+)=e^{-2\chi(0)}$, hence
only through the constant $\mathsf w_nH_n(0^+)$, which cancels in $\omega=D_W^2/W_0$ (see~\eqref{eq:DW-factor}).

\begin{lemma}[Corner modulus of the seam scalar]\label{lem:corner}
Let $0<p<1$ and $\tau_0\in(0,1)$. There are $C=C(p,\tau_0)$ and $n_0$ such that for $n\ge n_0$:
\begin{enumerate}
\item[(i)] $|\chi_\pm(\zeta)|\le C|\zeta|$ for $\zeta\in\Sseam_\pm$, $|\zeta|\le1$; $|\chi_\pm(\zeta)|\le C(1+\log|\zeta|)$ for
  $|\zeta|\ge1$; and $|\chi_\pm'(\zeta)|\le C$ for $|\zeta|\le2$.
\item[(ii)] For $z$ with $\mathrm{dist}(z,\Sseam_+\cup\Sseam_-)\ge\tau_0|z|$:
  $|\chi(z)-\chi(0)|\le C|z|(1+\log\frac1{|z|})$ and $|\chi'(z)|\le C(1+\log\frac1{|z|})$ if $|z|\le1$, and
  $|\chi(z)-\chi(0)|\le C(1+\log|z|)^2$ if $|z|\ge1$. Here
  $\chi(0):=\frac1{2\pi i}\bigl[\int_{\Sseam_+}\chi_+(\zeta)\frac{\diff\zeta}\zeta-\int_{\Sseam_-}\chi_-(\zeta)\frac{\diff\zeta}\zeta\bigr]$.
\item[(iii)] For each fixed $n$, $\chi$ is bounded on $U(0,1)\setminus(\Sseam_+\cup\Sseam_-)$, uniformly up to the seams.
\end{enumerate}
\end{lemma}

\begin{proof}
(i) On $\Sseam_+$, $t=-i|t|$ with $|t|=n|\zeta|^{-1/p}$. If $|\zeta|\le1$, then $|t|\ge n$; by~\eqref{eq:binet-bound} with
$\phi=\pi/2$, $|j(t)|\le\frac1{6|t|}\le|\zeta|^{1/p}/n\le|\zeta|$; $|\Log(1-e^{-2\pi|t|})|\le2e^{-2\pi n|\zeta|^{-1/p}}\le C|\zeta|$;
$|\varpi\Log(1+\zeta)|\le C|\zeta|$; $|\Log(1-\zeta/(2\pi n))|\le C|\zeta|/n$. If $|\zeta|\ge1$ and $|t|\ge1$, then
$|j(t)|\le\frac16$; if $|t|<1$, then~\eqref{eq:binet-small} (with a bounded remainder on $\{|t|\le1,\re t\ge0\}$) and
$|\log|t||\le\frac1p\log|\zeta|$ give $|j(t)|\le C(1+\log|\zeta|)$. Next $|\Log(1-e^{-2\pi|t|})|\le C(1+|\log|t||)\le
C(1+\log|\zeta|)$, since $1-e^{-y}\ge(1-e^{-2\pi})y/(2\pi)$ for $0<y\le2\pi$. Finally $|1-\zeta/(2\pi n)|$ lies between
$\sin\thseam$ and $1+|\zeta|$, because $\zeta/(2\pi n)$ lies on the ray of angle $\thseam\in(0,\pi/2)$. For the derivative on
$|\zeta|\le2$ (where $|t|\ge n2^{-1/p}\ge1$): $\frac{\diff}{\diff\zeta}j(t)=-j'(t)t/(p\zeta)=O(|t|^{-1}|\zeta|^{-1})=
O(|\zeta|^{1/p-1}/n)$, which is bounded since $1/p>1$; the derivative of $\Log(1-e^{-2\pi it})$ is
$O(|t||\zeta|^{-1}e^{-2\pi|t|})$, also bounded; the other two terms have bounded derivatives. $\Sseam_-$ follows by
conjugation.

(ii) For $z\notin\Sseam_\pm$,
$\chi(z)-\chi(0)=\frac z{2\pi i}\sum_\pm(\pm1)\int_{\Sseam_\pm}\frac{\chi_\pm(\zeta)\diff\zeta}{\zeta(\zeta-z)}$ and
$\chi'(z)=\frac1{2\pi i}\sum_\pm(\pm1)\int_{\Sseam_\pm}\frac{\chi_\pm(\zeta)\diff\zeta}{(\zeta-z)^2}$. Under the hypothesis,
$|\zeta-z|\ge\tau_0|z|$, and $|\zeta-z|\ge|\zeta|/2$ if $|\zeta|\ge2|z|$; hence $|\zeta-z|\ge\frac{\tau_0}3(|\zeta|+|z|)$.
With~(i), for $|z|\le1$,
\[
  |\chi(z)-\chi(0)|\le C|z|\Bigl[\int_0^1\frac{\diff\rho}{\rho+|z|}+\int_1^\infty\frac{1+\log\rho}{\rho^2}\diff\rho\Bigr]\le
  C|z|\bigl(1+\log\tfrac1{|z|}\bigr),
\]
and $|\chi'(z)|\le C[\int_0^1\rho(\rho+|z|)^{-2}\diff\rho+\int_1^\infty(1+\log\rho)\rho^{-2}\diff\rho]\le C(1+\log\frac1{|z|})$.
For $|z|\ge1$ the same splitting gives
$|z|\int_0^1\frac{\diff\rho}{|z|}+|z|\int_1^\infty\frac{(1+\log\rho)\diff\rho}{\rho(\rho+|z|)}\le C(1+\log|z|)^2$.

(iii) Let $z\in U(0,1)$ be close to $\Sseam_+$, with nearest point $z^*\in\Sseam_+$, $r:=|z^*|$ (other $z$ are covered
by~(ii)). Split $\int_{\Sseam_+}\chi_+(\zeta)\frac{\diff\zeta}{\zeta-z}$ into the part $|\zeta-z^*|\le r/2$ and the rest. On
the first part write $\chi_+(\zeta)=\chi_+(z^*)+O(|\zeta-z^*|)$ by~(i): the constant contributes $\chi_+(z^*)$ times a
logarithm whose modulus is at most $\pi+\log C$, because the two end points $z^*\pm\frac r2e^{i\thseam}$ are at comparable
distances from $z$; since $|\chi_+(z^*)|\le Cr$, this is bounded, and so is the Lipschitz part. On the rest,
$|\zeta-z|\ge(|\zeta|+r)/C$, and the contribution is at most $C\int_{\Sseam_+}|\chi_+(\zeta)|(|\zeta|+r)^{-1}|\diff\zeta|$, which
is finite for fixed $n$ because $\chi_+(\zeta)=O(|\zeta|)$ at $0$ and $\chi_+(\zeta)/\zeta$ is integrable at $\infty$
(Lemma~\ref{lem:seam}(ii)). The integral over $\Sseam_-$ is bounded in the same way, since
$|\zeta-z|\ge(|\zeta|+|z|)/C$ for $\zeta\in\Sseam_-$.
\end{proof}
